% Introduction Section - AEM Version
% Molecular Engineering of Semi-Transparent Non-Fullerene Acceptors for Quantum-Enhanced Agrivoltaics

\section{Introduction}\label{sec:Introduction}

Semi-transparent organic photovoltaics (OPVs) based on non-fullerene acceptors (NFAs) have emerged as promising candidates for dual-use energy applications, with recent power conversion efficiencies exceeding \SI{18}{\percent}~\cite{Liu2024_AEM_NFA,Wang2023_AEM_Stability}. The polymer donor PM6 paired with Y6-BO derivatives has demonstrated particular promise, achieving balanced charge transport and tunable absorption spectra ideal for spectral engineering applications~\cite{Zhang2024_AEM_ML}. However, optimizing these materials for agrivoltaic deployment requires simultaneous optimization of photovoltaic performance and spectral transmission properties---a multi-objective challenge that remains unresolved.

Current semi-transparent OPV designs prioritize either electrical efficiency or optical transparency, treating the transmitted spectrum as a secondary consideration~\cite{Cui2021,Wu2024}. This classical approach overlooks the quantum mechanical nature of photosynthetic energy transfer, where electronic coherences coupled to vibrational modes facilitate efficient exciton transport~\cite{Engel2007,Panitchayangkoon2010}. Biological light-harvesting complexes operate in an intermediate coupling regime where structured environmental fluctuations enhance rather than suppress quantum transport phenomena~\cite{Ishizaki2009,Collini2010}.

The Fenna-Matthews-Olsen (FMO) complex serves as a benchmark system for studying these quantum effects, with each monomer containing \numrange{7}{8} bacteriochlorophyll-a molecules that funnel energy through coherent pathways~\cite{Fenna1974,Adolphs2006}. Crucially, the energy transfer efficiency depends sensitively on the spectral composition of incident light, with specific vibronic resonances at \SIlist{750;820}{\nano\meter} enhancing coherence lifetimes by \SIrange{20}{50}{\percent}~\cite{Mohseni2014,Hildner2013}. This spectral dependence presents a materials design opportunity: OPV transmission spectra can be engineered to target these quantum mechanical resonances.

Recent advances in NFA molecular engineering enable precise spectral tuning through side-chain modification and backbone functionalization~\cite{Li2020,Firdaus2019}. Power conversion efficiencies above \SI{18}{\percent} have been achieved in semi-transparent configurations, yet systematic design rules for co-optimizing electrical and biological performance remain absent from the literature. The convergence of tunable OPV spectra and quantum-enhanced photosynthesis motivates our spectral bath engineering approach.

We introduce a materials design paradigm wherein the OPV transmission function $T(\omega)$ is strategically engineered to shape the effective spectral density experienced by photosynthetic complexes. Unlike classical spectral optimization that maximizes total photon flux, quantum spectral bath engineering targets specific vibronic transitions to unlock coherence-assisted transport mechanisms. Our AI-driven discovery pipeline integrates equivariant graph neural networks (E(n)-GNN) with Process Tensor HOPS simulations to screen \num{2847} donor-acceptor combinations, identifying Pareto-optimal candidates satisfying electrical, biological, and sustainability criteria simultaneously.

Operational stability represents a critical requirement for commercial deployment. The optimized PM6/Y6-BO derivative system demonstrates T$_{80}$ lifetimes exceeding \SI{18000}{\hour} under continuous illumination, with thermal cycling stability ($<$\SI{0.17}{\percent} degradation per cycle) projecting \num{20}+ year operational lifetimes~\cite{Wang2023_AEM_Stability}. Techno-economic modeling indicates levelized energy costs of \SIrange{0.04}{0.06}{\USD\per\kWh}, competitive with conventional silicon photovoltaics while enabling agricultural revenue diversification.

This work establishes five key contributions to materials design for dual-use energy applications:
\begin{enumerate}
    \item Demonstration of \SI{18.8}{\percent} PCE in semi-transparent OPV with quantum-enhanced biological compatibility (\SI{80}{\percent} photosynthetic efficiency maintained);
    \item Identification of vibronic resonance mechanism extending coherence lifetimes by \SIrange{20}{50}{\percent} through non-Markovian dynamics;
    \item AI-driven materials discovery pipeline screening \num{2847} candidates via E(n)-GNN prediction of quantum dynamical properties;
    \item Quantitative design rules from multi-objective Pareto optimization, balancing PCE (\SIrange{15.4}{22.1}{\percent}) with electron transport enhancement (\SIrange{12}{33}{\percent});
    \item Validation across nine climate zones confirming geographic robustness and commercial viability (LCOE \SIrange{0.04}{0.06}{\USD\per\kWh}).
\end{enumerate}

These findings bridge molecular-level materials engineering with device-level performance, providing a scalable methodology for next-generation semi-transparent OPV design.

The remainder of this paper is organized as follows. \Cref{sec:Theory} describes the open quantum system framework and PT-HOPS/SBD methodology. \Cref{sec:OPV_Materials} details the molecular design of semi-transparent OPVs and their optoelectronic characterization. \Cref{sec:Results} presents the quantum kinetics results and multi-objective Pareto optimization. \Cref{sec:Discussion} provides an economic and scalability analysis, while \Cref{sec:Conclusion} summarizes the key findings and future outlook.

\begin{figure*}[t]
\centering
\includegraphics[width=.7\textwidth]{Conceptual_Overview.png}
\caption{\textbf{Conceptual overview of quantum-enhanced agrivoltaic system.} \textbf{(a)} Schematic illustration of semi-transparent OPV panels integrated above agricultural crops, enabling simultaneous energy generation and photosynthesis. \textbf{(b)} Molecular structures of PM6 donor and Y6-BO acceptor materials with normalized absorption (solid lines) and transmission (shaded regions) spectra. Transmission windows at \SIlist{750;820}{\nano\meter} (dark green shading) coincide with FMO vibronic resonances. \textbf{(c)} Energy level diagram showing HOMO/LUMO alignment at the donor-acceptor interface with $\Delta E_{\mathrm{LUMO}} = \SI{0.75}{\electronvolt}$ and $\Delta E_{\mathrm{HOMO}} = \SI{0.35}{\electronvolt}$, facilitating efficient exciton dissociation and charge extraction.}
\label{fig:conceptual_overview}
\end{figure*}
